\documentclass{article}

\title{Snapshot Objectives\\Final Project}
\author{Elvis Montiel, David Montiel, Steven Mayorquin, Vincent Lozano}
\date{Snapshot 2}

\begin{document}
\maketitle

\section{Snapshot 1: Start}

The goal of Snapshot 1 is to establish the beginning of the simulated Want2Remember project. This includes creating the GitHub repository structure, adding the required practice simulation disclaimer, creating the initial README, creating the first SDD and SRS documents, creating the Docker Compose service plan, setting up Jira tasks and bugs, and adding initial TestRail reports.

\section{Snapshot 2: Checkpoint 1}

The goal of Snapshot 2 is to simulate the first major feature of the project: memory creation and memory templates. This snapshot builds on the initial project setup from Snapshot 1 by adding requirements, design details, user instructions, Jira tasks, bugs, and TestRail test cases for the memory creation workflow.

\subsection{Snapshot 2 Objectives}

The objectives for Snapshot 2 are to define the memory creation workflow, document memory template options, update the Software Design Document, update the Software Requirements Specification, update the Design Specification, update the User Manual, create new Jira tasks and bugs, and create a Snapshot 2 TestRail report.

\subsection{Snapshot 1 Reflection}

Snapshot 1 completed the initial repository structure, project disclaimer, README, initial SDD, initial SRS, design spec, user manual, Docker Compose service plan, Jira setup, and TestRail setup. The main remaining item from Snapshot 1 is the workflow diagram, which will be added or finalized in a later snapshot.

\section{Snapshot 3: Checkpoint 2}

The planned goal of Snapshot 3 is to simulate adding contacts and search features. This snapshot will update the project documents, Jira tasks, TestRail reports, and workflow diagram.

\section{Snapshot 4: Final Checkpoint}

The planned goal of Snapshot 4 is to finalize the repository, complete documentation, add final TestRail reports, update the workflow diagram, and describe future work.

\end{document}